\documentclass[14pt]{extarticle}
\usepackage{fontspec}
\usepackage{polyglossia}
\usepackage{amsmath}
\usepackage{amssymb}
\usepackage{xcolor}
\usepackage{fancyhdr}
\usepackage{graphicx}
\usepackage{listings}
\usepackage[most]{tcolorbox}
% \usepackage{setspace}
\usepackage{pifont}
\usepackage{enumitem}
\usepackage{titlesec}
\usepackage[bottom]{footmisc}
\usepackage{titling}
\usepackage{minted}
\usepackage{etoolbox}
\usepackage{array}
\usepackage{longtable} % page-breakable tables (must load before polyglossia/bidi RTL setup)
\usepackage{extsizes}
\usepackage{float}
\usepackage{booktabs}
\usepackage{hyperref}
\usepackage[onehalfspacing]{setspace}
\newif\ifwhiteblack
\ifdefined\setwhiteblack
  \setwhiteblack
\fi

\ifwhiteblack
  \usepackage[a4paper, landscape, top=1.25cm, bottom=1.25cm, left=1.25cm, right=1.25cm, headheight=1.5cm, headsep=0.5cm]{geometry}
\else
  \usepackage[a4paper, top=2.7cm, bottom=2.7cm, left=2.5cm, right=2.5cm, headheight=1.5cm, headsep=0.5cm]{geometry}
\fi
\usepackage{tikz}
\usepackage{datetime2}

\usetikzlibrary{automata, positioning, arrows.meta, arrows}
\providecommand{\LTcaptype}{table}

\newfontfamily\emoji{Segoe UI Emoji}

\pagestyle{fancy}

\setmainlanguage[numerals=western]{arabic}
\setotherlanguage{english}
\newfontfamily\arabicfont[Script=Arabic]{Calibri}
\newfontfamily\arabicfonttt[Script=Arabic]{Courier New}



\ifwhiteblack
\lstset{
  language=[Sharp]C,
  numbers=left,
  stepnumber=1,
  numbersep=8pt,
  frame=single,
  basicstyle=\ttfamily\small,
  keywordstyle=\bfseries,
  stringstyle=\ttfamily,
  commentstyle=\itshape
}
\else
\lstset{
  language=[Sharp]C,
  numbers=left,
  stepnumber=1,
  numbersep=8pt,
  frame=single,
  basicstyle=\ttfamily\small,
  keywordstyle=\color{blue},
  stringstyle=\color{red},
  commentstyle=\color{green!50!black}
}
\fi

\newif\ifdetailed
\ifdefined\setdetailed
  \setdetailed
\fi

\newif\ifwithsols
\ifdefined\setwithsols
  \setwithsols
\fi

\ifwhiteblack
  \AtBeginDocument{
    \LARGE
    \usemintedstyle{bw}
  }
\fi

% Default minted look: compact code that wraps long lines.
% Not applied in black-and-white landscape mode, which already uses \LARGE.
\ifwhiteblack\else
  \setminted{fontsize=\small, breaklines=true}
\fi

% unified tcolorboxes for articles
\ifwhiteblack
\tcbset{colback=white, colframe=black, fonttitle=\bfseries, boxrule=1pt}
\newtcolorbox{boxDef}[1][]{title={{\emoji📘} تعريف\ifx\\#1\\\else ~#1\fi :}}
\newtcolorbox{boxExercise}[1][]{title={{\emoji🧩} تمرين\ifx\\#1\\\else ~#1\fi :}}
\newtcolorbox{boxExample}[1][]{title={{\emoji📝} مثال\ifx\\#1\\\else ~#1\fi :}}
\newtcolorbox{boxNote}[1][]{title={{\emoji✨} ملاحظة\ifx\\#1\\\else ~#1\fi :}}
\newtcolorbox{boxAttention}[1][]{title={{\emoji🔔} تنبيه\ifx\\#1\\\else ~#1\fi :}}
\newtcolorbox{boxWarning}[1][]{title={{\emoji⚡} ملاحظة هامة\ifx\\#1\\\else ~#1\fi :}}
\newtcolorbox{boxSolution}[1][]{title={{\emoji✅} حل\ifx\\#1\\\else ~#1\fi :}}
\newtcolorbox{boxSymbol}[1][]{title={{\emoji🔣} رمز\ifx\\#1\\\else ~#1\fi :}}
\newtcolorbox{boxHint}[1][]{title={{\emoji💡} تلميح\ifx\\#1\\\else ~#1\fi :}}
\else
\tcbset{colback=white, colframe=black, fonttitle=\bfseries, boxrule=0.8pt}
\newtcolorbox{boxDef}[1][]{colback=blue!5!white,colframe=blue!75!black,
  title={{\emoji📘} تعريف\ifx\\#1\\\else ~#1\fi :}}
\newtcolorbox{boxExercise}[1][]{colback=cyan!5!white,colframe=cyan!70!black,
  title={{\emoji🧩} تمرين\ifx\\#1\\\else ~#1\fi :}}
\newtcolorbox{boxExample}[1][]{colback=yellow!5!white,colframe=orange!90!black,
  title={{\emoji📝} مثال\ifx\\#1\\\else ~#1\fi :}}
\newtcolorbox{boxNote}[1][]{colback=gray!10!white,colframe=black,
  title={{\emoji✨} ملاحظة\ifx\\#1\\\else ~#1\fi :}}
\newtcolorbox{boxAttention}[1][]{colback=magenta!10!white,colframe=magenta!80!black,
  title={{\emoji🔔} تنبيه\ifx\\#1\\\else ~#1\fi :}}
\newtcolorbox{boxWarning}[1][]{colback=red!5!white,colframe=red!75!black,
  title={{\emoji⚡} ملاحظة هامة\ifx\\#1\\\else ~#1\fi :}}
\newtcolorbox{boxSolution}[1][]{colback=green!5!white,colframe=green!60!black,
  title={{\emoji✅} حل\ifx\\#1\\\else ~#1\fi :}}
\newtcolorbox{boxSymbol}[1][]{colback=purple!5!white,colframe=purple!70!black,
  title={{\emoji🔣} رمز\ifx\\#1\\\else ~#1\fi :}}
\newtcolorbox{boxHint}[1][]{colback=teal!5!white,colframe=teal!60!black,
  title={{\emoji💡} تلميح\ifx\\#1\\\else ~#1\fi :}}
\fi


\ifwhiteblack
\tcbset{simplecode/.style={ colback=white, colframe=black, boxrule=1pt, arc=2pt, left=4pt,right=4pt,top=4pt,bottom=4pt}}
\else
\tcbset{simplecode/.style={ colback=gray!5, colframe=black!50, boxrule=0.4pt, arc=2pt, left=4pt,right=4pt,top=4pt,bottom=4pt}}
\fi
% Compact code boxes: tighter line spacing (the preamble uses onehalfspacing) and small vertical gaps.
% Must come AFTER the simplecode definition above.
\tcbset{simplecode/.append style={fontupper=\setstretch{0.95}, top=4pt, bottom=4pt, before skip=6pt, after skip=6pt}}
\newenvironment{boxCode}{\begin{tcolorbox}[simplecode]}{\end{tcolorbox}}

% Code cell used inside tables (keeps code left-to-right inside RTL tables).
% Use inside a fixed-width p{...} column: \begin{pycell}\begin{minted}...\end{minted}\end{pycell}
\newtcolorbox{pycell}{frame empty,colback=white,boxsep=0pt,left=2pt,right=2pt,top=2pt,bottom=2pt,width=\linewidth}

\newcolumntype{C}[1]{>{\centering\arraybackslash}p{#1}}

\makeatletter
\DTMnewdatestyle{dotted}{%
  \renewcommand{\DTMdisplaydate}[4]{\two@digits{##3}.\two@digits{##2}.##1}%
  \renewcommand{\DTMDisplaydate}{\DTMdisplaydate}%
}
\makeatother
\DTMsetdatestyle{dotted}

% redefine spaces after titles
\makeatletter
\renewcommand{\@maketitle}{%
  \begin{center}
    {\huge \bfseries \@title \par}%
    \vskip 0.2em % space between title and author
    {\large \@author \par}%
    % \vskip 0.2em % space between author and date
    % {\normalsize \@date \par}%
    {\small \hfill \begin{english}\DTMtoday\end{english} \par}%
  \end{center}
}
\makeatother

\fancyhf{} % clear default
\fancypagestyle{plain}{
  \fancyhf{}
  \fancyhead[L]{مدرسة التسامح الشاملة}
  % \fancyhead[L]{\includegraphics[height=1cm]{../../../images/logoTasamoh.png}}
  \fancyhead[R]{الأستاذ محمود اغبارية}
  \fancyfoot[C]{\thepage}
}

\fancyhead[L]{مدرسة التسامح الشاملة}
\fancyhead[R]{الأستاذ محمود اغبارية}
\fancyfoot[C]{\thepage}

\setcounter{tocdepth}{3} % only section subsection and subsubsection in TOC

\newcommand{\importpdfpage}[2]{%
    \thispagestyle{empty}
    \newgeometry{left=0cm,right=0cm,top=0cm,bottom=0cm}
    \includegraphics[page=#2,width=0.95\paperwidth,height=0.95\paperheight,keepaspectratio]{#1}
    \restoregeometry
}

\newcommand{\insertFullImg}[1]{%
    \noindent
    \makebox[\textwidth][c]{
        \includegraphics[width=0.9\paperwidth,keepaspectratio]{#1}%
    }%
}%


\ifwithsols
\title{حل ورقة تمرن $9$ للصف العاشر $10$ - الحلقات: الخانات، القواسم، والمتواليات}
\else
\title{ورقة تمرن $9$ للصف العاشر $10$ - الحلقات: الخانات، القواسم، والمتواليات}
\fi

\begin{document}

\maketitle
\thispagestyle{fancy}

% ==============================================================
% القسم الأول: أسئلة تحليل العدد إلى خاناته
% ==============================================================
\section*{القسم الأوّل: تحليل العدد إلى خاناته}

\begin{enumerate}[itemsep=2em]

% --- أ: استخراج خانات ثابتة (الآحاد والعشرات) ---
\item اكتب عملية خارجية تتلقى عددًا صحيحًا موجبًا $n$، وتطبع الأعداد من $1$ حتى $n$ التي تكون منزلة آحادها $5$ فقط.
\ifwithsols
\begin{boxSolution}
\begin{english}
\begin{minted}{csharp}
public static void PrintEndingWith5(int n)
{
    for (int i = 1; i <= n; i++)
    {
        if (i % 10 == 5)
        {
            Console.WriteLine(i);
        }
    }
}
\end{minted}
\end{english}
\end{boxSolution}
\fi
\filbreak

\item اكتب عملية خارجية تتلقى عددًا صحيحًا موجبًا $n$، وتعيد مجموع كل الأعداد بين $1$ و $n$ التي تكون منزلة آحادها $3$.
\ifwithsols
\begin{boxSolution}
\begin{english}
\begin{minted}{csharp}
public static int SumEnding3(int n)
{
    int sum = 0;
    for (int i = 1; i <= n; i++)
    {
        if (i % 10 == 3)
        {
            sum += i;
        }
    }
    return sum;
}
\end{minted}
\end{english}
\end{boxSolution}
\fi
\filbreak

\item اكتب برنامجًا يستقبل من المستخدم $20$ عددًا صحيحًا، ويطبع عدد الأعداد التي كانت منزلة آحادها أكبر من منزلة عشراتها.
\ifwithsols
\begin{boxSolution}
\begin{english}
\begin{minted}{csharp}
int cnt = 0;
for (int i = 0; i < 20; i++)
{
    int x = int.Parse(Console.ReadLine());
    x = Math.Abs(x);
    int ones = x % 10;
    int tens = (x / 10) % 10;
    if (ones > tens)
    {
        cnt++;
    }
}
Console.WriteLine(cnt);
\end{minted}
\end{english}
\end{boxSolution}
\fi
\filbreak

\item اكتب برنامجًا يستقبل من المستخدم أعدادًا صحيحة غير سالبة، ويعدّ الأعداد التي تكون منزلة العشرات فيها زوجية ويطبع عددها. \\
ينتهي الاستقبال عندما يُدخل المستخدم عددًا سالبًا.
\ifwithsols
\begin{boxSolution}[1]
\begin{english}
\begin{minted}{csharp}
int cnt = 0;
int n = int.Parse(Console.ReadLine());
while (n >= 0)
{
    int tens = (n / 10) % 10;
    if (tens % 2 == 0)
    {
        cnt++;
    }
    n = int.Parse(Console.ReadLine());
}
Console.WriteLine(cnt);
\end{minted}
\end{english}
\end{boxSolution}

\filbreak

\begin{boxSolution}[2: فكّر لماذا هو صحيح؟]
\begin{english}
\begin{minted}{csharp}
int cnt = 0;
int n = int.Parse(Console.ReadLine());
while (n >= 0)
{
    n /= 10;
    if (n % 2 == 0)
    {
        cnt++;
    }
    n = int.Parse(Console.ReadLine());
}
Console.WriteLine(cnt);
\end{minted}
\end{english}
\end{boxSolution}
\fi
\filbreak

% --- ب: تفكيك العدد بحلقة while (العدّ، الجمع، الضرب، الأكبر) ---
\item \label{item:digitCount} اكتب عملية خارجية باسم \textenglish{DigitsCount} تتلقى عددًا صحيحًا غير سالب، وتعيد عدد خاناته.
\ifwithsols
\begin{boxSolution}
\begin{english}
\begin{minted}{csharp}
public static int DigitsCount(int n)
{
    if (n == 0)
    {
        return 1;
    }
    int count = 0;
    while (n > 0)
    {
        count++;
        n /= 10;
    }
    return count;
}
\end{minted}
\end{english}
\end{boxSolution}
\fi
\filbreak

\item اكتب عملية خارجية باسم \textenglish{DigitsSum} تتلقى عددًا صحيحًا غير سالب، وتعيد مجموع خاناته.
\ifwithsols
\begin{boxSolution}
\begin{english}
\begin{minted}{csharp}
public static int DigitsSum(int n)
{
    int sum = 0;
    while (n > 0)
    {
        sum += (n % 10);
        n /= 10;
    }
    return sum;
}
\end{minted}
\end{english}
\end{boxSolution}
\fi
\filbreak

\item اكتب عملية خارجية تتلقى عددًا صحيحًا غير سالب، وتعيد حاصل ضرب خاناته. \\
\textbf{مثال:} إذا كان العدد $234$، تكون النتيجة $2 \times 3 \times 4 = 24$.
\ifwithsols
\begin{boxSolution}
\begin{english}
\begin{minted}{csharp}
public static int ProductOfDigits(int n)
{
    if (n == 0)
    {
        return 0;
    }
    int prod = 1;
    while (n != 0)
    {
        prod *= (n % 10);
        n /= 10;
    }
    return prod;
}
\end{minted}
\end{english}
\end{boxSolution}
\fi
\filbreak

\item اكتب عملية خارجية تتلقى عددًا صحيحًا موجبًا وتطبع: عدد خاناته، ومجموع خاناته، ومعدّل خاناته. \\
لا تستخدم العمليات من الأسئلة السابقة، واكتب حلقة \textenglish{while} واحدة فقط. \\
\textbf{مثال:} إذا تلقّت العملية العدد $762945$، فإنّها تطبع:
\begin{boxCode}
\begin{english}
\begin{minted}{text}
Digits count = 6
Digits sum = 33
Digits average = 5.5
\end{minted}
\end{english}
\end{boxCode}
\ifwithsols
\begin{boxSolution}
\begin{english}
\begin{minted}{csharp}
public static void PrintDigitsStats(int n)
{
    int count = 0, sum = 0;
    int m = n;
    while (m > 0)
    {
        sum += (m % 10);
        count++;
        m /= 10;
    }
    Console.WriteLine("Digits count = " + count);
    Console.WriteLine("Digits sum = " + sum);
    double avg = (double)sum / count;
    Console.WriteLine("Digits average = " + avg);
}
\end{minted}
\end{english}
\end{boxSolution}
\fi
\filbreak

\item اكتب عملية خارجية تتلقى عددًا صحيحًا غير سالب، وتطبع عدد الخانات الزوجية فيه وعدد الخانات الفردية فيه.
\ifwithsols
\begin{boxSolution}
\begin{english}
\begin{minted}{csharp}
public static void PrintEvenOddDigitsCount(int n)
{
    int even = 0, odd = 0;
    if (n == 0)
    {
        even = 1;
    }
    while (n > 0)
    {
        int d = n % 10;
        if (d % 2 == 0)
        {
            even++;
        }
        else
        {
            odd++;
        }
        n /= 10;
    }
    Console.WriteLine("Even digits = " + even);
    Console.WriteLine("Odd digits = " + odd);
}
\end{minted}
\end{english}
\end{boxSolution}
\fi
\filbreak

\item اكتب برنامجًا يستقبل من المستخدم $100$ عدد صحيح موجب، وفي النهاية يطبع العدد الذي له أكبر مجموع خانات (يمكنك الاستعانة بالعملية \textenglish{DigitsSum} التي كتبتها في الأسئلة السابقة).
\ifwithsols
\begin{boxSolution}
\begin{english}
\begin{minted}{csharp}
int bestNum = 0, bestSum = -1;
for (int i = 0; i < 100; i++)
{
    int x = int.Parse(Console.ReadLine());
    int s = DigitsSum(x);
    if (s > bestSum)
    {
        bestSum = s;
        bestNum = x;
    }
}
Console.WriteLine(bestNum);
\end{minted}
\end{english}
\end{boxSolution}
\fi
\filbreak

\item
\begin{enumerate}
    \item اكتب عملية خارجية باسم \textenglish{MaxDigit} تتلقى عددًا صحيحًا غير سالب وتعيد قيمة أكبر خانة فيه. \\
    \textbf{مثال:} إذا تلقّت العدد $23481$ فإنّها تعيد $8$.
    \item اكتب برنامجًا يستقبل من المستخدم $100$ عدد صحيح موجب، وفي النهاية يطبع العدد الذي ظهرت فيه أكبر خانة. إذا كان هناك أكثر من عدد ظهرت فيه أكبر خانة، يطبع البرنامج العدد الأكبر من بينها.
\end{enumerate}
\ifwithsols
\begin{boxSolution}[فرع أ]
\begin{english}
\begin{minted}{csharp}
public static int MaxDigit(int n)
{
    if (n == 0)
    {
        return 0;
    }
    int md = 0;
    while (n > 0)
    {
        md = Math.Max(md, n % 10);
        n /= 10;
    }
    return md;
}
\end{minted}
\end{english}
\end{boxSolution}

\filbreak

\begin{boxSolution}[فرع ب]
\begin{english}
\begin{minted}{csharp}
public static void Main(string[] args)
{
    int bestNum = 0, bestMax = -1;
    for (int i = 0; i < 100; i++)
    {
        int x = int.Parse(Console.ReadLine());
        int m = MaxDigit(x);
        if (m > bestMax)
        {
            bestMax = m;
            bestNum = x;
        }
        else if (m == bestMax && x > bestNum)
        {
            bestNum = x;
        }
    }
    Console.WriteLine(bestNum);
}
\end{minted}
\end{english}
\end{boxSolution}
\fi
\filbreak

% --- ج: البحث في الخانات وفحص خصائصها (شرط وجودي وشرط كلّي) ---
\item اكتب عملية خارجية باسم \textenglish{DigitExists} تتلقى عددًا صحيحًا غير سالب $n$ ورقمًا مكوّنًا من منزلة واحدة $d$ ($0 \le d \le 9$). \\
تعيد العملية \textenglish{true} إذا كان الرقم $d$ يظهر كإحدى خانات العدد $n$، وتعيد \textenglish{false} خلاف ذلك. \\
\textbf{مثال:} استدعاء العملية مع العدد $4729$ والرقم $7$ يعيد \textenglish{true}، بينما استدعاؤها مع العدد $4729$ والرقم $5$ يعيد \textenglish{false}.
\ifwithsols
\begin{boxSolution}
\begin{english}
\begin{minted}{csharp}
public static bool DigitExists(int n, int d)
{
    if (n == 0 && d == 0)
    {
        return true;
    }
    while (n > 0)
    {
        if (n % 10 == d)
        {
            return true;
        }
        n /= 10;
    }
    return false;
}
\end{minted}
\end{english}
\end{boxSolution}
\fi
\filbreak

\item اكتب عملية خارجية باسم \textenglish{AreDigitsAscending} تتلقى عددًا صحيحًا موجبًا $n$، وتعيد \textenglish{true} إذا كانت خانات العدد مرتّبة تصاعديًّا تمامًا من اليسار إلى اليمين (أي أنّ كل خانة أكبر تمامًا من الخانة التي على يسارها)، وتعيد \textenglish{false} خلاف ذلك. \\
\textbf{مثال:} للعدد $1359$ وللعدد $7$ تعيد العملية \textenglish{true}، بينما للعدد $1428$ وللعدد $1225$ تعيد \textenglish{false}.
\ifwithsols
\begin{boxSolution}
\begin{english}
\begin{minted}{csharp}
public static bool AreDigitsAscending(int n)
{
    while (n >= 10)
    {
        int rightDigit = n % 10;
        int leftDigit = (n / 10) % 10;
        if (rightDigit <= leftDigit)
        {
            return false;
        }
        n /= 10;
    }
    return true;
}
\end{minted}
\end{english}
\end{boxSolution}
\fi
\filbreak

% --- د: مجموع قوى الخانات والأعداد المميّزة ---
\item
\begin{enumerate}
    \item اكتب عملية خارجية باسم \textenglish{SumSquares} تتلقى عددًا صحيحًا موجبًا، وتعيد مجموع \textbf{تربيع} خاناته. \\
    \textbf{مثال:} إذا تلقّت العملية العدد $123$، فإنّها تعيد: $1^2 + 2^2 + 3^2 = 1 + 4 + 9 = 14$.
    \item في البرنامج الرئيسي، استدعِ العملية مع العدد $23$ واطبع النتيجة.
    \item اكتب مقطع برنامج في \textenglish{Main} يطبع جميع الأعداد بين $1$ و $10000$ التي تساوي مجموع تربيع خاناتها.
\end{enumerate}
\ifwithsols
\begin{boxSolution}
\begin{english}
\begin{minted}{csharp}
public static int SumSquares(int n)
{
    int sum = 0;
    while (n > 0)
    {
        int d = n % 10;
        sum += d * d;
        n /= 10;
    }
    return sum;
}

public static void Main(string[] args)
{
    Console.WriteLine(SumSquares(23));

    for (int i = 1; i <= 10000; i++)
    {
        if (SumSquares(i) == i)
        {
            Console.WriteLine(i);
        }
    }
}
\end{minted}
\end{english}
\end{boxSolution}
\fi
\filbreak

\item نعرّف «عددًا مميّزًا» بأنّه عدد موجب إذا رفعنا كل خانة من خاناته لقوة مساوية لعدد خاناته ثم جمعنا النتائج فإنّنا نحصل على العدد الأصلي نفسه. \\
\textbf{مثال:} العدد $153$ هو عدد مميّز، لأنّه مكوّن من $3$ خانات، و:
$1^3 + 5^3 + 3^3 = 1 + 125 + 27 = 153$.
\begin{enumerate}
    \item اكتب عملية خارجية باسم \textenglish{IsSpecial} تتلقى عددًا صحيحًا موجبًا، وتعيد \textenglish{true} إذا كان العدد مميّزًا، و \textenglish{false} خلاف ذلك. \\
    يمكنك الاستعانة بالعملية \textenglish{DigitsCount} من السؤال (\ref{item:digitCount}) لمعرفة عدد خانات العدد.
    \item في البرنامج الرئيسي، اطبع جميع الأعداد المميّزة من $1$ إلى $10000$ بالاستعانة بالعملية \textenglish{IsSpecial}.
\end{enumerate}
\ifwithsols
\begin{boxSolution}
\begin{english}
\begin{minted}{csharp}
public static bool IsSpecial(int n)
{
    int k = DigitsCount(n);
    int m = n, sum = 0;
    while (m > 0)
    {
        int d = m % 10;
        sum += (int)Math.Pow(d, k);
        m /= 10;
    }
    if (sum == n)
    {
        return true;
    }
    else
    {
        return false;
    }
}

public static void Main(string[] args)
{
    for (int i = 1; i <= 10000; i++)
    {
        if (IsSpecial(i))
        {
            Console.WriteLine(i);
        }
    }
}
\end{minted}
\end{english}
\end{boxSolution}
\fi
\filbreak

% --- هـ: بناء عدد، عكس عدد، والعدد المتماثل (Palindrome) ---
\item
\begin{enumerate}
    \item اكتب عملية خارجية باسم \textenglish{AppendDigit} تتلقى عددًا صحيحًا غير سالب، ورقمًا مكوّنًا من منزلة واحدة. \\
    العملية تلصق المنزلة في يمين العدد الأوّل وتعيد النتيجة. \\
    \textbf{مثال:} إذا تلقّت العملية العدد $12$ والمنزلة $5$ فإنّها تعيد $125$.
    \item اكتب برنامجًا يستقبل من المستخدم أرقامًا (كل منها مكوّن من منزلة واحدة) ويكوّن عددًا صحيحًا من هذه الأرقام ويطبعه. \\
    ينتهي استقبال الأرقام عندما يُدخل المستخدم عددًا سالبًا. \\
    \textbf{مثال:} إذا أدخل المستخدم الأرقام $5$ ثم $1$ ثم $2$ ثم $4$ ثم $-1$ فعلى البرنامج أن يطبع $5124$. \\
    عليك استخدام العملية \textenglish{AppendDigit} من الفرع السابق (لا حاجة للتحقق من صحة المدخلات).
\end{enumerate}
\ifwithsols
\begin{boxSolution}
\begin{english}
\begin{minted}{csharp}
public static int AppendDigit(int n, int d)
{
    return n * 10 + d;
}

public static void Main(string[] args)
{
    int res = 0;
    int d = int.Parse(Console.ReadLine());
    while (d >= 0)
    {
        res = AppendDigit(res, d);
        d = int.Parse(Console.ReadLine());
    }
    Console.WriteLine(res);
}
\end{minted}
\end{english}
\end{boxSolution}
\fi
\filbreak

\item
\begin{enumerate}
    \item اكتب عملية خارجية باسم \textenglish{ReverseNumber} تتلقى عددًا صحيحًا غير سالب $n$، وتعيد العدد الناتج من عكس ترتيب خاناته. \\
    \textbf{مثال:} إذا تلقّت العملية العدد $1234$ فإنّها تعيد $4321$.
    \item نسمّي العدد «متماثلاً» (بوليندروم) إذا كان يقرأ من اليمين إلى اليسار كما يقرأ من اليسار إلى اليمين (مثل: $12321$، $4554$، $7$). \\
    اكتب عملية خارجية باسم \textenglish{IsPalindrome} تتلقى عددًا صحيحًا غير سالب، وتعيد \textenglish{true} إذا كان العدد متماثلاً، و \textenglish{false} خلاف ذلك (استعن بالعملية \textenglish{ReverseNumber}).
\end{enumerate}
\ifwithsols
\begin{boxSolution}
\begin{english}
\begin{minted}{csharp}
// a
public static int ReverseNumber(int n)
{
    int rev = 0;
    while (n > 0)
    {
        rev = rev * 10 + (n % 10);
        n /= 10;
    }
    return rev;
}

// b
public static bool IsPalindrome(int n)
{
    if (n == ReverseNumber(n))
    {
        return true;
    }
    else
    {
        return false;
    }
}
\end{minted}
\end{english}
\end{boxSolution}
\fi
\filbreak

\item \label{item:appendDigitToLeft}
اكتب عملية خارجية باسم \textenglish{AppendDigitToLeft} تتلقى عددًا صحيحًا موجبًا، ورقمًا مكوّنًا من منزلة واحدة. \\
العملية تلصق المنزلة بالعدد الأوّل \textbf{من اليسار} وتعيد النتيجة. \\
\textbf{مثال:} إذا تلقّت العملية العدد $12$ والمنزلة $5$ فإنّها تعيد $512$.
\begin{boxHint}
يمكنك الاستعانة بالعملية \textenglish{DigitsCount} التي كتبتها في السؤال (\ref{item:digitCount}) لمعرفة عدد خانات العدد. \\
لكي نحصل على $512$ من $5$ و $12$ نضرب الـ $5$ بـ $100$ (وهو $10$ مرفوعًا لقوة عدد خانات الـ $12$) ثم نجمع الـ $12$.
\end{boxHint}
\ifwithsols
\begin{boxSolution}
\begin{english}
\begin{minted}{csharp}
public static int AppendDigitToLeft(int n, int d)
{
    int dc = DigitsCount(n);
    int p = (int)Math.Pow(10, dc);
    return d * p + n;
}
\end{minted}
\end{english}
\end{boxSolution}
\fi
\filbreak

\end{enumerate}

\subsection*{سؤال 3 امتحان 899371 سنة 2025}

\insertFullImg{../../../bagrut_questions/basics/loops_both_2025_899371_3.png}

\ifwithsols
\begin{boxSolution}[سؤال 3 امتحان 899371 سنة 2025]
\begin{enumerate}[label=\alph*.]
    \item .
\begin{boxCode}
\begin{english}
\begin{minted}{csharp}
private static int DigitsCount(int n)
{
    int cnt = 0;
    while (n > 0)
    {
        cnt++;
        n /= 10;
    }
    return cnt;
}
public static bool IsNarc(int num)
{
    int digitsCnt = DigitsCount(num);
    int sum = 0, digit;
    int temp = num;
    while (temp > 0)
    {
        digit = temp % 10;
        sum += (int)Math.Pow(digit, digitsCnt);
        temp /= 10;
    }
    return sum == num;
}
\end{minted}
\end{english}
\end{boxCode}

\item .
\begin{boxCode}
\begin{english}
\begin{minted}{csharp}
public static void TheNarc(int n)
{
    for (int i = 1; i <= n; i++)
    {
        if (IsNarc(i))
            Console.WriteLine(i);
    }
}
\end{minted}
\end{english}
\end{boxCode}
\end{enumerate}
\end{boxSolution}
\clearpage
\fi

\importpdfpage{../../../bagrut_questions/basics/loops_both_2024_899371_6.pdf}{1}

\ifwithsols

\paragraph*{حل سؤال 6 امتحان 899371 سنة 2024}:
\bigskip

% // TODO:

\fi


\clearpage
% ==============================================================
% القسم الثاني: أسئلة تحليل العدد إلى عوامله (قواسمه)
% ==============================================================
\section*{القسم الثاني: تحليل العدد إلى عوامله (قواسمه) والأعداد الأوّلية}

\begin{enumerate}[resume, itemsep=2em]

% --- أ: إيجاد القواسم، عددها، مجموعها، والأعداد الكاملة والصديقة ---
\item اكتب برنامجًا يستقبل من المستخدم عددًا صحيحًا موجبًا، ويطبع جميع عوامله (قواسمه)، أي كل الأعداد الموجبة التي يقبل هذا العدد القسمة عليها بدون باقٍ.
\ifwithsols
\begin{boxSolution}
\begin{english}
\begin{minted}{csharp}
int n = int.Parse(Console.ReadLine());
for (int d = 1; d <= n; d++)
{
    if (n % d == 0)
    {
        Console.WriteLine(d);
    }
}
\end{minted}
\end{english}
\end{boxSolution}
\fi
\filbreak

\item اكتب عملية خارجية باسم \textenglish{CountDivisors} تتلقى عددًا صحيحًا موجبًا $n$، وتعيد عدد قواسمه الموجبة (بما فيها الـ $1$ والعدد $n$ نفسه). \\
\textbf{مثال:} إذا تلقّت العملية العدد $12$ فإنّها تعيد $6$ (لأنّ قواسم الـ $12$ هي: $1, 2, 3, 4, 6, 12$).
\ifwithsols
\begin{boxSolution}
\begin{english}
\begin{minted}{csharp}
public static int CountDivisors(int n)
{
    int count = 0;
    for (int d = 1; d <= n; d++)
    {
        if (n % d == 0)
        {
            count++;
        }
    }
    return count;
}
\end{minted}
\end{english}
\end{boxSolution}
\fi
\filbreak

\item \label{item:divisors_sum}
اكتب عملية خارجية باسم \textenglish{DivisorsSum} تتلقى عددًا صحيحًا موجبًا وتعيد \textbf{مجموع} قواسمه \textbf{باستثناء نفسه}. \\
\textbf{مثال:} إذا تلقّت العدد $8$، فإنّها تعيد: $1 + 2 + 4 = 7$.
\ifwithsols
\begin{boxSolution}
\begin{english}
\begin{minted}{csharp}
public static int DivisorsSum(int n)
{
    int sum = 0;
    for (int d = 1; d < n; d++)
    {
        if (n % d == 0)
        {
            sum += d;
        }
    }
    return sum;
}
\end{minted}
\end{english}
\end{boxSolution}
\fi
\filbreak

\item \textbf{العدد الكامل:} هو عدد موجب يساوي مجموع قواسمه باستثناء نفسه. \\
\textbf{مثال:} العدد $6$ هو عدد كامل لأنّ $1 + 2 + 3 = 6$. \\
اكتب عملية خارجية باسم \textenglish{IsPerfectNumber} تتلقى عددًا صحيحًا موجبًا وتعيد \textenglish{true} إذا كان عددًا كاملاً، و \textenglish{false} وإلا (يمكنك الاستعانة بالعملية \textenglish{DivisorsSum} من السؤال \ref{item:divisors_sum}).
\ifwithsols
\begin{boxSolution}[1]
\begin{english}
\begin{minted}{csharp}
public static bool IsPerfectNumber(int n)
{
    int sum = 0;
    for (int i = 1; i < n; i++)
    {
        if (n % i == 0)
        {
            sum += i;
        }
    }
    return sum == n;
}
\end{minted}
\end{english}
\end{boxSolution}

\filbreak

\begin{boxSolution}[2: مع استعمال العملية \textenglish{DivisorsSum}]
\begin{english}
\begin{minted}{csharp}
public static bool IsPerfectNumber(int n)
{
    return DivisorsSum(n) == n;
}
\end{minted}
\end{english}
\end{boxSolution}
\fi
\filbreak

\item \textbf{عددان صديقان:} هما عددان موجبان يحقّقان أنّ مجموع قواسم الأوّل (باستثناء نفسه) يساوي العدد الثاني، ومجموع قواسم الثاني (باستثناء نفسه) يساوي العدد الأوّل. \\
\textbf{مثال:} العددان $220$ و $284$ هما عددان صديقان، لأنّ مجموع قواسم $220$ (بدون نفسه) يساوي $284$، ومجموع قواسم $284$ (بدون نفسه) يساوي $220$. \\
اكتب عملية خارجية باسم \textenglish{AreAmicable} تتلقى عددين صحيحين موجبين وتعيد \textenglish{true} إذا كانا عددين صديقين، و \textenglish{false} وإلا (يمكنك الاستعانة بالعملية \textenglish{DivisorsSum} من السؤال \ref{item:divisors_sum}).
\ifwithsols
\begin{boxSolution}
\begin{english}
\begin{minted}{csharp}
public static bool AreAmicable(int a, int b)
{
    if (DivisorsSum(a) == b && DivisorsSum(b) == a)
    {
        return true;
    }
    else
    {
        return false;
    }
}
\end{minted}
\end{english}
\end{boxSolution}
\fi
\filbreak

% --- ب: الأعداد الأوّلية وتطبيقاتها ---
\item \textbf{تذكير:} العدد الأوّلي هو عدد صحيح أكبر من $1$ لا يقبل القسمة إلا على نفسه وعلى $1$.
\begin{enumerate}
    \item اكتب عملية خارجية باسم \textenglish{IsPrime} تتلقى عددًا صحيحًا وتعيد \textenglish{true} إذا كان هذا العدد أوّليًّا، وإلا تعيد \textenglish{false}.
    \item في البرنامج الرئيسي، استدعِ العملية \textenglish{IsPrime} مع العدد $17$ واطبع النتيجة، ثم مع العدد $20$ واطبع النتيجة.
    \item اكتب عملية خارجية أخرى باسم \textenglish{PrintPrimes} تتلقى عددًا صحيحًا $n$ وتطبع كل الأعداد الأوّلية من $2$ حتى $n$ (بالاستعانة بالعملية \textenglish{IsPrime}). \\
    \textbf{مثال:} إذا تلقّت العملية العدد $20$ فإنّها تطبع: $2, 3, 5, 7, 11, 13, 17, 19$.
    \item في البرنامج الرئيسي، استدعِ العملية \textenglish{PrintPrimes} مع العدد $100$.
\end{enumerate}
\ifwithsols
\begin{boxSolution}[فرع أ و ج]
\begin{english}
\begin{minted}{csharp}
// a
public static bool IsPrime(int n)
{
    if (n < 2)
    {
        return false;
    }
    for (int i = 2; i < n; i++)
    {
        if (n % i == 0)
        {
            return false;
        }
    }
    return true;
}

// c
public static void PrintPrimes(int n)
{
    for (int i = 2; i <= n; i++)
    {
        if (IsPrime(i))
        {
            Console.WriteLine(i);
        }
    }
}
\end{minted}
\end{english}
\end{boxSolution}

\filbreak

\begin{boxSolution}[فرع ب و د]
\begin{english}
\begin{minted}{csharp}
public static void Main(string[] args)
{
    // b
    Console.WriteLine(IsPrime(17));
    Console.WriteLine(IsPrime(20));

    // d
    PrintPrimes(100);
}
\end{minted}
\end{english}
\end{boxSolution}
\fi
\filbreak

\item \textbf{الأعداد الأوّلية التوأم:} هما عددان أوّليان الفرق بينهما يساوي $2$ (مثل: $(3, 5)$، $(5, 7)$، $(11, 13)$، $(17, 19)$). \\
اكتب عملية خارجية باسم \textenglish{PrintTwinPrimes} تتلقى عددًا صحيحًا موجبًا $n$، وتطبع جميع أزواج الأعداد الأوّلية التوأم التي لا تتجاوز $n$ (استعن بالعملية \textenglish{IsPrime}).
\ifwithsols
\begin{boxSolution}
\begin{english}
\begin{minted}{csharp}
public static void PrintTwinPrimes(int n)
{
    for (int i = 2; i <= n - 2; i++)
    {
        if (IsPrime(i) && IsPrime(i + 2))
        {
            Console.WriteLine(i + " , " + (i + 2));
        }
    }
}
\end{minted}
\end{english}
\end{boxSolution}
\fi
\filbreak

\item اكتب عملية خارجية تتلقى عددًا صحيحًا موجبًا، وتطبع كل عوامله (قواسمه) \textbf{الأوّلية} (أي كل الأعداد الأوّلية التي يقبل هذا العدد القسمة عليها بدون باقٍ). \\
يمكنك الاستعانة بالعملية \textenglish{IsPrime} التي كتبتها في الأسئلة السابقة.
\ifwithsols
\begin{boxSolution}
\begin{english}
\begin{minted}{csharp}
public static void PrintPrimeFactors(int n)
{
    for (int i = 2; i <= n; i++)
    {
        if (n % i == 0 && IsPrime(i))
        {
            Console.WriteLine(i);
        }
    }
}
\end{minted}
\end{english}
\end{boxSolution}
\fi
\filbreak

\item اكتب عملية خارجية باسم \textenglish{IsPrimeSquare} تتلقى عددًا صحيحًا موجبًا وتعيد \textenglish{true} إذا كان العدد الذي تلقّته عبارة عن تربيع عدد أوّلي، وإلا تعيد \textenglish{false}. \\
\textbf{مثال:} تعيد العملية \textenglish{true} إذا تلقّت $25$ (لأنّ $25 = 5^2$ والعدد $5$ أوّلي)، وتعيد \textenglish{false} إذا تلقّت $30$ أو $36$.
\ifwithsols
\begin{boxSolution}[1]
\begin{english}
\begin{minted}{csharp}
public static bool IsPrimeSquare(int n)
{
    for (int i = 2; i <= Math.Sqrt(n); i++)
    {
        if (i * i == n && IsPrime(i))
        {
            return true;
        }
    }
    return false;
}
\end{minted}
\end{english}
\end{boxSolution}

\filbreak

\begin{boxSolution}[2]
\begin{english}
\begin{minted}{csharp}
public static bool IsPrimeSquare(int n)
{
    int sq = (int)Math.Sqrt(n);
    if (sq * sq == n && IsPrime(sq))
    {
        return true;
    }
    else
    {
        return false;
    }
}
\end{minted}
\end{english}
\end{boxSolution}
\fi
\filbreak

% --- ج: القاسم المشترك الأكبر (GCD) والمضاعف المشترك الأصغر (LCM) ---
\item اكتب عملية خارجية باسم \textenglish{GCD} تتلقى عددين صحيحين موجبين $a$ و $b$، وتعيد القاسم المشترك الأكبر لهما (أكبر عدد يقبل كلا العددين القسمة عليه بدون باقٍ). \\
\textbf{مثال:} القاسم المشترك الأكبر للعددين $12$ و $15$ هو $3$. \\
\textbf{ملاحظة:} القاسم المشترك الأكبر يكون بالضرورة أصغر من أو مساويًا لـ \textenglish{Math.Min(a, b)}.
\ifwithsols
\begin{boxSolution}
\begin{english}
\begin{minted}{csharp}
public static int GCD(int a, int b)
{
    int min = Math.Min(a, b);
    for (int i = min; i >= 1; i--)
    {
        if (a % i == 0 && b % i == 0)
        {
            return i;
        }
    }
    return 1;
}
\end{minted}
\end{english}
\end{boxSolution}
\fi
\filbreak

\item اكتب عملية خارجية باسم \textenglish{GCDThreeNumbers} تتلقى ثلاثة أعداد صحيحة موجبة $a$ و $b$ و $c$، وتعيد القاسم المشترك الأكبر للأعداد الثلاثة. \\
\textbf{مثال:} القاسم المشترك الأكبر للأعداد $12$ و $18$ و $24$ هو $6$. \\
\textbf{ملاحظة:} يمكنك الاستعانة بالعملية \textenglish{GCD} التي كتبتها في السؤال السابق.
\ifwithsols
\begin{boxSolution}[1]
\begin{english}
\begin{minted}{csharp}
public static int GCDThreeNumbers(int a, int b, int c)
{
    int min = Math.Min(c, Math.Min(a, b));
    for (int i = min; i >= 1; i--)
    {
        if (a % i == 0 && b % i == 0 && c % i == 0)
        {
            return i;
        }
    }
    return 1;
}
\end{minted}
\end{english}
\end{boxSolution}

\filbreak

\begin{boxSolution}[2]
\begin{english}
\begin{minted}{csharp}
public static int GCDThreeNumbers(int a, int b, int c)
{
    int gcdAB = GCD(a, b);
    return GCD(gcdAB, c);
}
\end{minted}
\end{english}
\end{boxSolution}
\fi
\filbreak

\item \textbf{سؤال للتفكير:} معطاة العملية \textenglish{GCDThreeNumbers} التي كتبتها في السؤال السابق. كيف يمكنك باستخدامها فقط كتابة العملية \textenglish{GCD} التي تحسب القاسم المشترك الأكبر لعددين صحيحين $a$ و $b$؟
\ifwithsols
\begin{boxSolution}
\begin{english}
\begin{minted}{csharp}
public static int GCD(int a, int b)
{
    return GCDThreeNumbers(a, a, b);
}
\end{minted}
\end{english}
\end{boxSolution}
\fi
\filbreak

\item \textbf{عددان أوّليان نسبيًّا:} هما عددان القاسم المشترك الأكبر لهما يساوي $1$. \\
\textbf{مثال:} العددان $15$ و $28$ أوّليان نسبيًّا لأنّ القاسم المشترك الأكبر لهما يساوي $1$. \\
اكتب عملية خارجية باسم \textenglish{AreRelativelyPrime} تتلقى عددين صحيحين موجبين وتعيد \textenglish{true} إذا كانا أوّليين نسبيًّا، و \textenglish{false} وإلا (يمكنك الاستعانة بالعملية \textenglish{GCD}).
\ifwithsols
\begin{boxSolution}
\begin{english}
\begin{minted}{csharp}
public static bool AreRelativelyPrime(int a, int b)
{
    return GCD(a, b) == 1;
}
\end{minted}
\end{english}
\end{boxSolution}
\fi
\filbreak

\item اكتب عملية خارجية باسم \textenglish{LCM} تتلقى عددين صحيحين موجبين $a$ و $b$، وتعيد المضاعف المشترك الأصغر لهما (أصغر عدد موجب يقبل القسمة على العددين معًا بدون باقٍ). \\
\textbf{مثال:} المضاعف المشترك الأصغر للعددين $6$ و $8$ هو $24$، لأنّه أصغر عدد موجب يقبل القسمة على $6$ وعلى $8$ معًا. \\
\textbf{ملاحظة:} المضاعف المشترك الأصغر يكون بالضرورة أكبر من أو مساويًا لـ \textenglish{Math.Max(a, b)} وأصغر من أو مساويًا لـ $a \cdot b$.
\ifwithsols
\begin{boxSolution}
\begin{english}
\begin{minted}{csharp}
public static int LCM(int a, int b)
{
    int max = Math.Max(a, b);
    for (int i = max; i < (a * b); i++)
    {
        if (i % a == 0 && i % b == 0)
        {
            return i;
        }
    }
    return a * b;
}
\end{minted}
\end{english}
\end{boxSolution}
\fi
\filbreak

\end{enumerate}

\clearpage
% ==============================================================
% القسم الثالث: أسئلة المتواليات والسلاسل
% ==============================================================
\section*{القسم الثالث: أسئلة المتواليات والسلاسل}

\begin{enumerate}[resume, itemsep=2em]

\item اكتب برنامجًا يطبع أوّل $20$ حدًّا من المتوالية: $3, 6, 9, 12, \dots$
\ifwithsols
\begin{boxSolution}
\begin{english}
\begin{minted}{csharp}
for (int i = 1; i <= 20; i++)
{
    Console.WriteLine(3 * i);
}
\end{minted}
\end{english}
\end{boxSolution}
\fi
\filbreak

\item اكتب عملية خارجية تتلقى عددًا صحيحًا موجبًا $n$، وتطبع أوّل $n$ حدود من المتوالية التالية:
\[ 1,\; 2,\; 4,\; 7,\; 11,\; 16,\; \ldots \]
(نحصل على حدود هذه المتوالية بأن نبدأ بالعدد $1$، ثم نضيف $1$، ثم نضيف $2$، ثم نضيف $3$، ثم نضيف $4$، وهكذا.)
\ifwithsols
\begin{boxSolution}
\begin{english}
\begin{minted}{csharp}
public static void PrintSequence(int n)
{
    int val = 1;
    for (int i = 1; i <= n; i++)
    {
        Console.WriteLine(val);
        val += i;
    }
}
\end{minted}
\end{english}
\end{boxSolution}
\fi
\filbreak

\item اكتب عملية خارجية تتلقى عددًا صحيحًا موجبًا $n$، وتحسب وتعيد مجموع السلسلة التالية:
\[ \text{Sum} = 1 + \frac{1}{2} + \frac{1}{3} + \dots + \frac{1}{n} \]
\ifdetailed
\begin{boxExample}
\begin{itemize}[nosep]
    \item إذا تلقّت $1$ فإنّها تعيد $1$.
    \item إذا تلقّت $2$ فإنّها تعيد: $1 + \frac{1}{2} = 1.5$.
    \item إذا تلقّت $3$ فإنّها تعيد: $1 + \frac{1}{2} + \frac{1}{3} = 1.8333\dots$
    \item إذا تلقّت $10$ فإنّها تعيد: $1 + \frac{1}{2} + \frac{1}{3} + \ldots + \frac{1}{10} = 2.9289\dots$
\end{itemize}
\end{boxExample}
\fi
\ifwithsols
\begin{boxSolution}
\begin{english}
\begin{minted}{csharp}
public static double SeriesSum(int n)
{
    double sum = 0;
    for (int i = 1; i <= n; i++)
    {
        sum += 1.0 / i;
    }
    return sum;
}
\end{minted}
\end{english}
\end{boxSolution}
\fi
\filbreak

\item اكتب عملية خارجية تتلقى عددًا صحيحًا موجبًا $n$، وتحسب وتعيد مجموع السلسلة التالية:
\[ \text{Sum} = \frac{1}{2} + \frac{2}{3} + \frac{3}{4} + \dots + \frac{n}{n+1} \]
\ifwithsols
\begin{boxSolution}
\begin{english}
\begin{minted}{csharp}
public static double FractionsSeriesSum(int n)
{
    double sum = 0;
    for (int i = 1; i <= n; i++)
    {
        sum += (double)i / (i + 1);
    }
    return sum;
}
\end{minted}
\end{english}
\end{boxSolution}
\fi
\filbreak

\item اكتب عملية خارجية باسم \textenglish{MinTermsToExceed} تتلقى عددًا عشريًّا موجبًا $target$، وتجد وتعيد أصغر عدد حدود $n$ يجب جمعها من السلسلة:
\[ 1 + \frac{1}{2} + \frac{1}{3} + \dots + \frac{1}{n} \]
لكي يصبح المجموع أكبر تمامًا من $target$. \\
\textbf{مثال:} إذا تلقّت العملية العدد $1.7$ فإنّها تعيد $3$ (لأنّ $1 + \frac{1}{2} = 1.5 \le 1.7$، بينما $1 + \frac{1}{2} + \frac{1}{3} \approx 1.833 > 1.7$).
\ifwithsols
\begin{boxSolution}
\begin{english}
\begin{minted}{csharp}
public static int MinTermsToExceed(double target)
{
    double sum = 0;
    int n = 0;
    while (sum <= target)
    {
        n++;
        sum += 1.0 / n;
    }
    return n;
}
\end{minted}
\end{english}
\end{boxSolution}
\fi
\filbreak

\item متوالية فيبوناتشي معرّفة بالشكل التالي:
\begin{english}
\begin{align*}
    a_1 &= 1 \\
    a_2 &= 1 \\
    a_n &= a_{n-1} + a_{n-2} \quad \text{for } n > 2
\end{align*}
\end{english}
مثلاً، أوّل $7$ حدود من متوالية فيبوناتشي هي: $1, 1, 2, 3, 5, 8, 13$.
\begin{enumerate}
    \item اكتب عملية خارجية باسم \textenglish{Fib} تتلقى عددًا صحيحًا موجبًا $n$ وتعيد الحد الـ $n$ من متوالية فيبوناتشي (مثلاً: إذا تلقّت العدد $3$ تعيد $2$، وإذا تلقّت $6$ تعيد $8$).
    \item اكتب عملية خارجية باسم \textenglish{PrintFirstNFibs} تتلقى عددًا صحيحًا موجبًا $n$، وتطبع أوّل $n$ حدود من متوالية فيبوناتشي بالاستعانة بالعملية \textenglish{Fib} من البند (أ). \\
    مثلاً، إذا تلقّت العملية العدد $8$، تطبع الأعداد: $1, 1, 2, 3, 5, 8, 13, 21$ (كل عدد في سطر).
    \item اكتب برنامجًا رئيسيًّا يستقبل عددًا صحيحًا موجبًا من المستخدم، ويستدعي العملية \textenglish{PrintFirstNFibs} مع العدد المدخل.
\end{enumerate}
\ifwithsols
\begin{boxSolution}[فرع أ]
\begin{english}
\begin{minted}{csharp}
public static int Fib(int n)
{
    if (n <= 2)
    {
        return 1;
    }
    int a = 1, b = 1;
    for (int i = 3; i <= n; i++)
    {
        int c = a + b;
        a = b;
        b = c;
    }
    return b;
}
\end{minted}
\end{english}
\end{boxSolution}

\filbreak

\begin{boxSolution}[فرع ب و ج]
\begin{english}
\begin{minted}{csharp}
public static void PrintFirstNFibs(int n)
{
    for (int i = 1; i <= n; i++)
    {
        Console.WriteLine(Fib(i));
    }
}

public static void Main(string[] args)
{
    int n = int.Parse(Console.ReadLine());
    PrintFirstNFibs(n);
}
\end{minted}
\end{english}
\end{boxSolution}
\fi
\filbreak

\end{enumerate}

\vspace{1cm}
\begin{flushleft}
أرجو لكم وقتًا ممتعًا.

الأستاذ محمود اغبارية.
\end{flushleft}

\end{document}
